\section{Integración DevOps}\label{sec:devops}

\subsection{Entorno de desarrollo reproducible}

Todo el equipo trabaja en el mismo \textbf{contenedor de desarrollo} (\emph{devcontainer}), una pila de Docker
Compose con dos servicios: el terminal y DynamoDB Local en memoria. La imagen trae las versiones que usa la CI (Node~24,
Python~3.12 con \code{uv}, Terraform~1.16, AWS CLI, \code{gh}, \code{just}, \code{crane}),
Claude Code y, opcionalmente, otros agentes de terminal; nueve volúmenes con nombre conservan
dependencias, sesiones de los agentes y la caché de SSO entre reconstrucciones. El
guion \emph{post-create} ejecuta \code{just setup} (ficheros \code{.env}, dependencias de Python y Node)
e instala Chromium para Playwright. El acceso a AWS es solo por \textbf{IAM Identity Center}
(\code{just aws-login}): ni las personas ni los agentes del contenedor manejan claves de larga
duración (Fig.~\ref{fig:devcontainer}).

\begin{figure}[h]
\centering
\resizebox{\columnwidth}{!}{\begin{tikzpicture}[x=1cm,y=1cm]
\node[draw=kgray,rounded corners=4pt,dashed,minimum width=6.35cm,minimum height=2.75cm] (host) at (3.05,0.05) {};
\node[font=\tiny\bfseries,color=kgray,anchor=north west] at (-0.1,1.42) {Docker Compose (\code{.devcontainer})};
\node[box,font=\tiny,text width=3.45cm,minimum height=1.45cm,align=left] (dc) at (1.85,0.4) {\textbf{devcontainer} (Ubuntu)\\Node 24 · Python 3.12 + uv\\Terraform · AWS CLI · gh · just\\crane · Chromium (Playwright)\\Claude Code (+ Codex, Gemini)};
\node[data,font=\tiny,text width=1.75cm,minimum height=1.45cm] (dyn) at (5.0,0.4) {\textbf{DynamoDB Local}\\en memoria};
\node[box,font=\tiny,text width=5.95cm,minimum height=0.45cm,fill=gray!8] (vol) at (3.05,-0.85) {9 volúmenes: \code{.venv}, \code{node\_modules}, sesiones de agentes, caché SSO};
\node[aws,font=\tiny,text width=1.3cm,minimum height=1.45cm] (sso) at (7.25,0.4) {\textbf{IAM Identity Center} (SSO)};
\node[box,font=\tiny,text width=7.9cm,minimum height=0.45cm] (run) at (3.85,1.95) {\code{just dev-core} :8000 · \code{just dev-ai} :8001 · \code{just dev-frontend} :3000 (recarga en caliente)};
\draw[arr] (dc) -- (dyn);
\draw[arr] (dc.north) -- (dc.north |- run.south);
\draw[darr] (sso.south) |- node[lbl,pos=0.75,below]{credenciales temporales (\code{just aws-login})} (3.3,-1.35) -- (3.3,-1.08);
\draw[darr] (vol.north -| 3.3,0) -- (dc.south -| 3.3,0);
\end{tikzpicture}
}
\caption{Entorno de desarrollo: el mismo contenedor para personas y agentes, con las herramientas de la CI y credenciales de AWS temporales.}\label{fig:devcontainer}
\end{figure}

\textbf{Una sola interfaz: \code{just}.} Personas, agentes y CI ejecutan las mismas recetas
(\code{just lint}, \code{just test}, \code{just contracts}, \code{just corpus}, \code{just index},
\code{just stack-up}\ldots). Un único repositorio (\emph{monorepo}) con un espacio de trabajo de \code{uv} (y un único
fichero de bloqueo de dependencias, \emph{lockfile}) agrupa los servicios, la biblioteca común y
las herramientas.

\subsection{Contenedores}\label{sec:contenedores}

Un único \code{Dockerfile} multietapa genera las dos imágenes del backend (argumento
\code{SERVICE}): una etapa instala las dependencias con \code{uv} a partir del \emph{lockfile}, sin
paquetes de desarrollo, y la final parte de \code{python:3.12-slim}, se ejecuta con un usuario sin
privilegios e incluye una comprobación de salud. La misma imagen funciona en Docker Compose y en
Lambda gracias a AWS Lambda Web Adapter, una extensión que traduce cada invocación a una petición
HTTP local y que fuera de Lambda no hace nada. Hay dos pilas locales: \code{just docker-up}
(solo backend: Nginx, ambos servicios y DynamoDB Local) y \code{just stack-up}, que añade el sitio
estático y reproduce el despliegue con un único origen en el puerto 8080; Nginx replica las reglas
de CloudFront, incluido no almacenar en búfer el flujo SSE. En cada propuesta de cambio que toca el backend, la CI construye ambas
imágenes y comprueba que la aplicación importa, que el adaptador está presente y que el servicio de
IA responde a la comprobación de salud dentro del contenedor.

\subsection{Integración y despliegue continuos}\label{sec:cicd}

\begin{figure*}[t]
\centering
\resizebox{\linewidth}{!}{\begin{tikzpicture}[x=1cm,y=1cm,
  a/.style={ai,text width=2.05cm,minimum height=0.8cm},
  c/.style={box,text width=2.05cm,minimum height=0.8cm},
  h/.style={human,text width=2.05cm,minimum height=0.8cm},
  w/.style={aws,text width=2.05cm,minimum height=0.8cm}]
\node[c] (lin) at (0,0) {\textbf{Tarea} en Linear\\y rama};
\node[a] (brief) at (2.55,0) {\textbf{Auditoría y \emph{brief}}\\modelo más capaz};
\node[a] (impl) at (5.1,0) {\textbf{Implementa}\\Opus en \emph{worktree}};
\node[c] (pr) at (7.65,0) {\textbf{Propuesta de cambio}\\enlazada a la tarea};
\node[a,text width=2.4cm] (rev) at (10.35,0) {\textbf{3 revisores} Sonnet\\corrección · contratos y pruebas · UX y accesibilidad};
\node[a] (fix) at (13.05,0) {\textbf{Corrige} y vigila la CI (Opus)};
\node[h] (merge) at (15.6,0) {\textbf{\emph{Merge}}\\(\emph{squash})};
\node[c,text width=7.2cm,minimum height=0.9cm,align=left] (ci) at (7.65,-1.55) {\textbf{Comprobaciones de la CI} (por rutas): ruff + pyright · pytest por paquete · ESLint + Vitest · Playwright (estático y pila Compose) · deriva de contratos OpenAPI$\rightarrow$TS · imagen Docker + prueba de humo · \code{terraform validate} · \code{docs-check}};
\node[w] (fe) at (15.6,-1.55) {\textbf{Frontend}: \code{s3 sync} + invalidación (auto)};
\node[c] (img) at (13.05,-2.95) {\textbf{Imágenes} en GHCR (automático)};
\node[h] (dep) at (15.6,-2.95) {\textbf{Despliegue del backend}: copia a ECR, \code{plan}, \code{apply}};
\draw[arr] (lin)--(brief); \draw[arr] (brief)--(impl); \draw[arr] (impl)--(pr); \draw[arr] (pr)--(rev);
\draw[arr] (rev)--(fix); \draw[arr] (fix)--(merge);
\draw[arr] (pr.south) -- (ci.north -| pr.south);
\draw[darr] (ci.east) -| node[lbl,pos=0.25,above]{estado de la CI} (fix.south);
\draw[arr] (merge) -- (fe);
\draw[arr] (merge.south west) -- (img.north);
\draw[arr] (img) -- (dep);
\node[lbl,anchor=north] at (2.55,-0.55) {\emph{decisiones cerradas,}\\\emph{pruebas y documentación}};
\end{tikzpicture}}
\caption{De la tarea de Linear a producción. En morado, agentes de IA; en verde, aprobación de una persona (el \emph{merge} la requiere cuando el sistema de permisos lo exige); en naranja, despliegues en AWS. La CI se ejecuta en cada propuesta de cambio y bloquea la integración si falla.}\label{fig:cicd}
\end{figure*}

La CI tiene 13 trabajos filtrados por rutas: un cambio solo en el frontend no
ejecuta las pruebas de Python. Además del análisis estático (\emph{lint}), la comprobación de
tipos y las pruebas, verifica tres cosas propias de
este proyecto: (1) \textbf{deriva de contratos}: si un esquema Pydantic cambia sin regenerar el
OpenAPI y los tipos TypeScript del frontend, la CI falla; (2) la \textbf{pila completa} (Nginx,
ambos servicios y DynamoDB Local en Compose) se levanta y se recorre con Playwright con sesión
iniciada; y (3) la \textbf{documentación}: ficheros obligatorios, índice de decisiones de arquitectura, enlaces y que
toda receta \code{just} citada en las instrucciones de los agentes exista. Hay 1.092 pruebas en
el backend (665 en \code{ai\_api} y 191 en el constructor de corpus), unas 900 de Vitest en el
frontend y unas 75 de Playwright.


\begin{table}[h]
\centering\scriptsize
\caption{Pruebas automáticas por paquete (funciones o casos de prueba; los parametrizados cuentan una vez).}\label{tab:tests}
\begin{tabularx}{\columnwidth}{@{}L{2.2cm}rY@{}}
\toprule
\textbf{Paquete} & \textbf{Pruebas} & \textbf{Qué cubren} \\
\midrule
\code{ai\_api} & 665 & Agente, herramientas, recuperación, \emph{guardrails}, SSE, trazas (con dobles de modelo y de base vectorial) \\
\code{city\_corpus} & 191 & Limpieza, fusión, troceo, validación y puerta, con \emph{fixtures} sin red \\
\code{core\_api} & 175 & Viajes, usuarios, acceso y cuota sobre DynamoDB simulado (moto) \\
\code{travel\_common} & 40 & JWT, errores, registros de acceso \\
\code{vector\_store\_bench} & 21 & Métricas y búsqueda del experimento \\
Frontend (Vitest) & $\approx$900 & Componentes, reductor del planificador, i18n completo \\
Frontend (Playwright) & $\approx$75 & Flujos en desarrollo, sitio estático y pila completa \\
\bottomrule
\end{tabularx}
\end{table}

El despliegue separa construir de promover: cada \emph{merge} que toca el backend construye imágenes en GHCR (el registro de
contenedores de GitHub); el despliegue del backend (manual) las copia a ECR (el de AWS) con \code{crane}, fija el \emph{digest} en
Terraform y aplica el plan solo si se pide. El frontend se publica automáticamente. Es una
decisión deliberada: el frontend practica despliegue continuo, porque publicarlo no cambia
infraestructura; el backend practica \emph{entrega} continua, porque cada \emph{merge} deja una
imagen lista para promover, pero aplicarla modifica recursos de pago y la decide una persona
después de revisar el plan de Terraform, la misma puerta humana que fija el proceso de entrega. GitHub Actions
se autentica en AWS por \textbf{OIDC} con un rol restringido al entorno \code{aws}. Dependabot
actualiza dependencias semanalmente y \code{just release} asigna la versión y etiqueta cada publicación.

\subsection{Infraestructura como código}

La infraestructura de AWS está en Terraform (\code{infra/aws}): CloudFront, S3, Route~53,
ACM, API Gateway, Lambda, Cognito, DynamoDB, S3 Vectors, Bedrock Guardrails, ECR, IAM
y los cuadros de mando de CloudWatch. Quedan fuera, a propósito, el WAF de CloudFront (opcional,
asociado desde la consola e ignorado de forma explícita por Terraform para no deshacerlo; llevarlo
a código es parte de la arquitectura objetivo, §\ref{sec:prod}) y el modo \code{ENHANCED} del
índice vectorial, que se fija con la CLI porque el proveedor aún no lo admite. El estado de Terraform se guarda en un \emph{bucket} de S3 cifrado
y versionado con bloqueo nativo, creado por una configuración raíz de arranque (\emph{bootstrap}) que también da de alta el proveedor
OIDC de GitHub y los roles de CI. Se conserva además una variante anterior para Google Cloud
(Cloud Run y Cloud SQL) que la CI valida sintácticamente; no está al día con DynamoDB ni con la
base vectorial, y se mantiene como punto de partida para una futura portabilidad.

\subsection{Entrega con agentes de IA}\label{sec:entrega}

Desde el 21 de septiembre, el proceso por defecto para todo cambio de más de unas pocas líneas tiene cuatro roles
(Fig.~\ref{fig:cicd}), orquestado por un \emph{workflow} de Claude Code
(\code{kyrian-wave.js}):
\begin{enumerate}
  \item \textbf{Auditoría y \emph{brief}}: el modelo más capaz lee el código, \emph{toma} las decisiones
        y escribe un \emph{brief} (especificación para el agente) con el resultado esperado, las
        decisiones cerradas, los ficheros, las pruebas, la documentación y los criterios de
        aceptación. La tarea de Linear guarda el resultado y la lista de aceptación.
  \item \textbf{Implementación}: un agente Opus trabaja en su propia copia de trabajo aislada
        (\emph{git worktree}), ejecuta las comprobaciones y abre la propuesta de cambio.
  \item \textbf{Revisión}: tres agentes Sonnet, en paralelo y cada uno con un enfoque (corrección;
        contratos y pruebas; UX, textos y accesibilidad), devuelven hallazgos estructurados con
        fichero, línea y severidad.
  \item \textbf{Corrección y entrega}: Opus aplica los hallazgos confirmados (o justifica por qué
        no), vigila la CI durante un máximo de dos rondas y nunca integra: lo hace la sesión
        principal con la CI en verde. La persona aprueba el \emph{merge} si el sistema de permisos
        lo exige, lanza \code{terraform apply} y prueba en producción.
\end{enumerate}
El principio es que el modelo más capaz decide y los más económicos implementan y revisan, y que
cada cambio llegue a una persona con un \emph{brief} escrito, tres revisiones y la CI en verde. Las
reglas se aplican con herramientas, no con buena voluntad: un \emph{hook} (comprobación automática
de Claude Code) impide crear ramas cuyo nombre no indique el tipo de cambio y la tarea de Linear
a la que responde, \code{AGENTS.md} es la fuente única de instrucciones para cualquier
agente (Claude Code, Codex, Gemini\ldots) y diez comandos encapsulan tareas repetidas
(empezar una tarea, entregarla con agentes, \code{/add-city}, \code{/check-site}\ldots). En cifras:
242 entregas integradas en la rama principal (92 de ellas desde la adopción del proceso); el
75\,\% de las entregas se generaron con Claude Code.

\subsection{Monitorización y registro}\label{sec:monitorizacion}

La observabilidad tiene tres niveles, todos sin servidores adicionales:
\begin{itemize}
  \item \textbf{Trazas de IA por turno} en DynamoDB, con el vocabulario de OpenInference
        y OpenTelemetry GenAI: tramos (\emph{spans}) de modelo (tokens, primer fragmento, esquema, ids
        elegidos y descartados), recuperación (consulta, filtros, distancias, escalón de
        ampliación), herramienta y cadena, más la línea temporal SSE y el coste. Se escriben antes
        de cerrar el flujo, caducan a los 90 días y se exploran en la consola \code{/admin}, que
        incluye un inspector de turnos con vista en cascada.
  \item \textbf{Registros (\emph{logs}) estructurados}: una línea \code{access\_log=} JSON por petición (ruta como
        plantilla, estado, duración, arranque en frío) sin datos personales, y eventos de negocio
        \code{business\_event=} seudonimizados al crear un viaje.
  \item \textbf{Cuadros de mando de CloudWatch} definidos en Terraform: servicio (Lambda, API
        Gateway, Cognito, Bedrock), infraestructura (DynamoDB, CloudFront) y registros de acceso con
        Logs Insights, más un panel de negocio en Grafana sobre los eventos.
\end{itemize}
La Tabla~\ref{tab:spans} muestra la anatomía de la traza del turno de la Tabla~\ref{tab:traza}:
basta abrirla en el inspector para saber qué buscó el agente, con qué filtros, qué documentos
devolvió la base vectorial (con su distancia), cuáles acabaron en el itinerario y cuánto costó.
Con ello la trazabilidad del dato es completa: \textbf{fuente} (URL y licencia en el documento)
$\rightarrow$ \textbf{documento} (\code{doc\_id} en Git) $\rightarrow$ \textbf{vector} (clave derivada
del id) $\rightarrow$ \textbf{búsqueda} (id y distancia en la traza) $\rightarrow$ \textbf{tarjeta}
del itinerario (mismo id, con su fuente visible para el viajero).
Las métricas concretas se describen en la §\ref{sec:metricas}.

\begin{table}[h]
\centering\scriptsize
\caption{\emph{Spans} de la traza de un turno de edición, en orden.}\label{tab:spans}
\begin{tabularx}{\columnwidth}{@{}L{2.35cm}L{1.1cm}Y@{}}
\toprule
\textbf{\emph{Span}} & \textbf{Tipo} & \textbf{Qué registra} \\
\midrule
\code{read\_turn}, \code{guard} & cadena & Idioma, ids ya usados; veredicto de Bedrock Guardrails \\
\code{agent:1} & modelo & Tokens, tiempo al primer fragmento, herramientas pedidas, motivo de parada \\
\code{tool:search\_places} & herram. & Argumentos; dentro, un \emph{span} de recuperación con consulta, filtros, $k$, escalón y distancias \\
\code{agent:2} & modelo & Tres llamadas de escritura \\
\code{tool:add\_activity} & herram. & Id, día y franja; búsqueda de foto \\
\code{agent:3} & modelo & Respuesta final (\code{end\_turn}) \\
\code{validate\_day}, \code{strip\_prices} & cadena & Avisos emitidos; importes eliminados \\
Resumen & -- & Tokens, coste, latencia, documentos recuperados y usados, ids descartados \\
\bottomrule
\end{tabularx}
\end{table}

\subsection{Reproducibilidad}

Todo resultado de esta memoria se puede regenerar desde el repositorio con las mismas recetas
(\code{index} y las evaluaciones requieren una sesión de AWS; la del agente, Bedrock o una clave de NVIDIA):

\noindent\colorbox{klight}{\parbox{\dimexpr\columnwidth-2\fboxsep}{\ttfamily\scriptsize
just setup \&\& just lint \&\& just test \hfill\rmfamily\emph{comprobaciones de la CI}\ttfamily\\
just corpus budapest \hfill\rmfamily\emph{corpus + puerta de calidad}\ttfamily\\
just index budapest \hfill\rmfamily\emph{embeddings e índice}\ttfamily\\
just eval-retrieval \hfill\rmfamily\emph{recall@k y MRR}\ttfamily\\
just eval-answers \hfill\rmfamily\emph{juez LLM}\ttfamily\\
just eval-conversations "" -{}-repeat 3 -{}-concurrency 2 \hfill\rmfamily\emph{agente}\ttfamily\\
just stack-up \hfill\rmfamily\emph{pila completa en :8080}}}
